\documentclass[a4paper,12pt]{article}
\usepackage[utf8]{inputenc}
\usepackage[T1]{fontenc}
\usepackage[ngerman]{babel}
\usepackage{geometry}
\usepackage{amsmath}
\usepackage{amssymb}
\usepackage{booktabs}
\geometry{margin=2.5cm}

\begin{document}

% ================================================================
\section*{Scalable Transformer for PDE Surrogate Modeling}

\textbf{Autoren:} Zijie Li, Dule Shu, Amir Barati Farimani \\
\textbf{Institution:} Carnegie Mellon University, Department of Mechanical Engineering \\
\textbf{Konferenz:} NeurIPS 2023 \\
\textbf{arXiv:} 2305.17560

% ================================================================
\subsection*{1\quad Problemstellung und Motivation}

Transformer-basierte Operatorlerner haben vielversprechende Ergebnisse für PDE-Surrogat-Modellierung
gezeigt. Die Standardvariante des skalierten Skalarprodukts besitzt jedoch quadratische Komplexität
$\mathcal{O}(N^2 d)$ in der Anzahl der Gitterpunkte $N$. Für 2D-Gitter mit $256 \times 256 = 65\,536$
Punkten oder 3D-Gitter erzeugt dies eine extrem große Aufmerksamkeitsmatrix, was den
Einsatz auf hochauflösenden Gittern praktisch unmöglich macht.

Sofmax-freie linearisierte Aufmerksamkeit (z.B.\ Galerkin-Transformer, Cao~2021) reduziert
die Komplexität auf $\mathcal{O}(Nd^2)$, bleibt aber quadratisch im Kanalparameter $d$
und kann bei tiefen gestapelten Schichten numerisch instabil werden.

Das Paper schlägt den \textbf{FactFormer} vor -- einen Transformer mit faktorisiertem
Kernintegral, der die Fluch der Dimensionalität gezielt durch eine axiale Faktorisierung
des Kernintegrals bekämpft.

% ================================================================
\subsection*{2\quad Methodik: Faktorisiertes Kernintegral}

\paragraph{Kernintegral-Perspektive der Aufmerksamkeit.}
Unter dem Kernel-Integralrahmen entspricht der softmax-freie Aufmerksamkeitsmechanismus
einer numerischen Quadratur:
\begin{equation}
    (\mathbf{z}_i)^l \approx \int_\Omega \kappa(x_i, \xi)\, v_l(\xi)\,\mathrm{d}\xi,
\end{equation}
wobei $\kappa(x, \xi) = \sum_l q_l(x) k_l(\xi)$ ein instanzabhängiger lernbarer Kern ist.
Für hochdimensionale Gebiete ist die Auswertung dieses vollen Kernintegrals über alle
$N$ Gitterpunktpaare rechnerisch prohibitiv.

\paragraph{Lernbare Projektion.}
Der erste Kernbaustein von FactFormer ist ein Satz lernbarer Integraloperatoren
$\{\mathcal{G}^{(1)}, \ldots, \mathcal{G}^{(n)}\}$, der die Eingabefunktion
$u : \mathbb{R}^n \to \mathbb{R}^d$ auf $n$ Funktionen mit
\emph{eindimensionalem} Definitionsbereich projiziert:
\begin{equation}
    \phi^{(m)}(x_i^{(m)}) = \mathcal{G}^{(m)}(u)(x_i^{(m)})
    \approx h^{(m)}\!\left(w \int_{\Omega_{-m}} \gamma^{(m)}\!\left(u(\xi_1, \ldots, x_i^{(m)}, \ldots, \xi_n)\right) \mathrm{d}\xi_{-m}\right),
\end{equation}
d.h.\ durch Mean-Pooling über alle Achsen außer der $m$-ten und anschließender nichtlinearer
Abbildung. In der Praxis ist $\gamma^{(m)}$ eine lineare Transformation, $h^{(m)}$ ein
dreischichtiges MLP.

\paragraph{Faktorisiertes Kernintegral.}
Das faktorisierte Kernintegral ersetzt das volle $n$-dimensionale Integral durch ein
Produkt von $n$ eindimensionalen Kernintegralen:
\begin{equation}
    z\!\left(x_{i_1}^{(1)}, \ldots, x_{i_n}^{(n)}\right)
    = \int_{\Omega_1} \kappa^{(1)}(x_{i_1}^{(1)}, \xi_1) \cdots
      \int_{\Omega_n} \kappa^{(n)}(x_{i_n}^{(n)}, \xi_n)\,
      v(\xi_1, \ldots, \xi_n)\,\mathrm{d}\xi_1 \cdots \mathrm{d}\xi_n.
\end{equation}
Jeder Achsenkern $\kappa^{(m)}$ hat die Komplexität $\mathcal{O}(S_m^2 d)$,
wobei $S_m$ die Gittergröße entlang der $m$-ten Achse ist. Die Gesamtkomplexität
beträgt $\mathcal{O}(S_1^2 d + S_2^2 d + \ldots + S_n^2 d)$, also
$\mathcal{O}(n S_m^2 d)$ statt $\mathcal{O}(N^2 d)$ für das vollständige Integral
mit $N = S_1 \cdots S_n$.

Als Positionskodierung wird \emph{Rotary Positional Encoding} (RoPE) eingesetzt,
das das Skalarprodukt im Kernintegral mit der relativen Position moduliert:
\begin{equation}
    Z = w\,\widetilde{Q}\widetilde{K}^\top V, \quad
    \text{wobei } \widetilde{Q}_i = g(\mathbf{q}_i, x_i) = \mathbf{q}_i\,\boldsymbol{\Theta}(x_i).
\end{equation}

% ================================================================
\subsection*{3\quad Experimente und Ergebnisse}

Die Modelle werden mit \emph{Latent Marching} (LM) und \emph{Autoregressive} (AR) Rollout
trainiert. FactFormer wird gegen FNO2D/3D, F-FNO, Dil-ResNet und Galerkin-Transformer verglichen.

\begin{center}
\begin{tabular}{lrrr}
\toprule
\textbf{Modell} & \textbf{2D KF avg $\omega$} & \textbf{3D Rauch $d$ avg} & \textbf{Inf.-Zeit (s)} \\
\midrule
FNO2D/3D & 0.2978 & 0.1344 & 0.73/3.19 \\
F-FNO & 0.2811 & 0.1038 & 0.81/2.75 \\
Dil-ResNet & 0.1655 & 0.0999 & 1.78/6.94 \\
FactFormer & 0.3017 & 0.0942 & 1.38/2.62 \\
\bottomrule
\end{tabular}
\end{center}

Für 2D-Kolmogorov-Strömung ($256\times256$) ist FactFormer mit Dil-ResNet vergleichbar
und deutlich besser als FNO. Für 3D-Rauchauftrieb ($64\times64\times64$) ist FactFormer
das effizienteste Modell (2.62\,s) bei kompetitiver Genauigkeit. Bei 3D-isotroper
Turbulenz zeigt Dil-ResNet kurzfristig bessere Leistung, FactFormer profitiert aber
von geringerem Speicherbedarf (4.6\,M vs.\ 509\,M Parameter bei FNO3D).

Ein wichtiges Ergebnis ist die \textbf{Auflösungsinvarianz}: Transformer-basierte
Modelle (FNO, FactFormer) zeigen nahezu konstante Genauigkeit über verschiedene
Trainings-Diskretisierungen, während Dil-ResNet (CNN) stark an die Trainingsgröße gebunden ist.

% ================================================================
\subsection*{4\quad Bezug zur Masterarbeit}

\paragraph{4.1\quad Skalierbarkeit als strukturelle Herausforderung.}
In der Masterarbeit werden Stokes-Strömungen auf einem fixen $256 \times 256$-Gitter
modelliert -- einer Auflösung, bei der standard quadratische Aufmerksamkeit
($N = 65\,536$ Punkte) praktisch nicht anwendbar ist. FactFormer demonstriert
einen Weg, wie Transformer auch auf solchen Gittern einsetzbar werden, indem
die axiale Faktorisierung die effektive Sequenzlänge pro Achse auf $S_m = 256$ reduziert.
Dies unterstreicht indirekt die Designentscheidung in der Masterarbeit, auf
konvolutionsbasierte Architekturen (U-Net, FNO) zu setzen, die auf regulären
Gittern besonders effizient sind.

\paragraph{4.2\quad Elliptische Probleme: stationäre Darcy-Strömung.}
FactFormer wird auch auf stationäre Darcy-Strömung (2D) getestet und erzielt dort
mit dem Galerkin-Transformer vergleichbare Ergebnisse. Da die Stokes-Gleichungen
ebenfalls ein elliptischer Operator sind, spricht dieser Befund dafür, dass
faktorisierte Aufmerksamkeit grundsätzlich für das in der Masterarbeit betrachtete
Problem geeignet wäre -- mit dem Vorteil einer besseren Skalierbarkeit als
das Galerkin-Modell bei hochauflösenden Gittern.

\paragraph{4.3\quad Auflösungsinvarianz als Vorteil gegenüber U-Net.}
Die beobachtete Auflösungsinvarianz von Transformer-Modellen gegenüber
CNN-basierten Methoden wie Dil-ResNet ist ein konzeptuell bedeutsames Ergebnis.
Das U-Net in der Masterarbeit ist ebenfalls ein CNN-basiertes Modell und somit
prinzipiell an die Trainingsgröße $256 \times 256$ gebunden. FactFormer liefert
damit ein Argument für die Überlegenheit spektral- oder aufmerksamkeitsbasierter
Methoden (zu denen auch FNO zählt) hinsichtlich Auflösungstransfer -- ein Aspekt,
der im Ausblick der Masterarbeit als zukünftige Forschungsrichtung relevant ist.

\paragraph{4.4\quad Faktorisierung vs.\ Spektral-Trunkierung: komplementäre Stärken.}
FNO trunkiert hohe Fourier-Moden, um globale Abhängigkeiten effizient zu erfassen,
büßt dabei aber Hochfrequenzgenauigkeit an Kristallgrenzen ein. FactFormer
faktorisiert das Kernintegral axial und behält prinzipiell Zugang zu allen
Ortsfrequenzen (sofern die Achsenkerne ausreichend ausdrucksstark sind).
Diese komplementäre Stärke motiviert FactFormer als potenzielle Architektur,
die die FNO-Schwäche an Kristallgrenzen überwindet, ohne die globale Reichweite
zu verlieren -- ein konkreter Forschungsausblick für Folgearbeiten zur Masterarbeit.

\paragraph{4.5\quad Abgrenzung.}
FactFormer wurde primär für zeitabhängige Strömungssimulationen entwickelt und
verwendet spezielle Trainingsstrategien (Latent Marching, Pushforward). Die
Masterarbeit behandelt ausschließlich stationäre Zustandsfelder, sodass diese
Trainingsaspekte nicht direkt übertragbar sind. Zudem variiert FactFormer die
geometrische Topologie nicht (feste Gitterstruktur ohne eingebettete Festkörper),
was die zentrale Generalisierungsfrage der Masterarbeit -- variierende Kristallzahlen
-- nicht adressiert.

\end{document}
